% =====================================================================
% Informe final · Estadística Multivariada · UCC Río Cuarto · 2.º sem. 2026
% Proyecto integrador sobre UNSW-NB15.
%
% Normas de la Etapa 5: A4 · márgenes 2,5 cm · fuente tipo Arial 11 pt ·
% interlineado 1,5 · justificado · número de página abajo a la derecha desde la
% portada · índices de secciones, tablas y gráficos · epígrafe en cada gráfico ·
% referencias en página nueva y en orden alfabético.
% Límite del grupo: el cuerpo (Introducción a Conclusiones) no supera 15 páginas.
%
% Compilar desde esta carpeta (informe/) con pdfLaTeX + BibTeX:
%   latexmk -pdf main.tex
% Las figuras se leen de ../resultados/figuras y los scripts de ../R y ../python,
% así que en Overleaf hay que subir el repositorio completo (o mantener esa estructura).
% =====================================================================
\documentclass[11pt,a4paper]{article}

\usepackage[utf8]{inputenc}
\usepackage[T1]{fontenc}
\usepackage[spanish,es-tabla,es-nodecimaldot]{babel}  % "Tabla"; coma decimal sin espacio
\usepackage{lmodern}

% --- Fuente: Helvetica/Nimbus Sans, métricamente equivalente a Arial ---
\usepackage[scaled=1]{helvet}
\renewcommand{\familydefault}{\sfdefault}

% --- Página ---
\usepackage[a4paper,margin=2.5cm]{geometry}

% --- Interlineado 1,5, sin espacio extra entre párrafos, justificado ---
\usepackage{setspace}
\onehalfspacing
\setlength{\parskip}{0pt}
\setlength{\parindent}{1cm}
\usepackage{microtype}

% --- Número de página abajo a la derecha (también en la portada) ---
\usepackage{fancyhdr}
\pagestyle{fancy}
\fancyhf{}
\fancyfoot[R]{\thepage}
\renewcommand{\headrulewidth}{0pt}
\fancypagestyle{plain}{\fancyhf{}\fancyfoot[R]{\thepage}\renewcommand{\headrulewidth}{0pt}}

% --- Títulos ---
\usepackage{titlesec}
\titleformat{\section}{\normalfont\large\bfseries}{\thesection.}{0.6em}{}
\titleformat{\subsection}{\normalfont\normalsize\bfseries}{\thesubsection.}{0.6em}{}
\titlespacing*{\section}{0pt}{10pt}{4pt}
\titlespacing*{\subsection}{0pt}{6pt}{2pt}

% --- Tablas, gráficos y epígrafes ---
\usepackage{graphicx}
\graphicspath{{../resultados/figuras/}}
\usepackage{booktabs}
\usepackage{tabularx}
\usepackage{array}
\usepackage{multirow}
\usepackage{float}
\usepackage[table]{xcolor}
\usepackage[font={small,stretch=1},labelfont=bf,labelsep=period,
            justification=justified,singlelinecheck=false,skip=4pt]{caption}
\captionsetup[table]{position=above}
\captionsetup[figure]{position=below}
% Gráficos "encuadrados" (la consigna pide gráficos encuadrados con su referencia)
\setlength{\fboxsep}{2pt}
\setlength{\fboxrule}{0.4pt}
\newcommand{\grafico}[2][\linewidth]{\fbox{\includegraphics[width=#1]{#2}}}
% Tablas compactas: interlineado simple y letra un punto menor dentro de tablas
\newenvironment{tablacompacta}{\begin{singlespace}\small}{\end{singlespace}}

% --- Referencias ---
\usepackage[authoryear,round]{natbib}
\setlength{\bibsep}{2pt}

% --- Código de los anexos ---
\usepackage{listings}
\definecolor{codigofondo}{HTML}{F7F7F5}
\definecolor{codigocomentario}{HTML}{5C7A5C}
\definecolor{codigoclave}{HTML}{1F4E79}
\lstset{
  basicstyle=\ttfamily\scriptsize\linespread{1.0}\selectfont,
  backgroundcolor=\color{codigofondo},
  commentstyle=\color{codigocomentario}\itshape,
  keywordstyle=\color{codigoclave}\bfseries,
  stringstyle=\color{black!70},
  numbers=left, numberstyle=\tiny\color{gray}, numbersep=6pt,
  frame=single, rulecolor=\color{black!25},
  breaklines=true, breakatwhitespace=false, columns=fullflexible, keepspaces=true,
  showstringspaces=false, tabsize=2, upquote=true,
  postbreak=\mbox{\textcolor{gray}{$\hookrightarrow$}\space},
  literate=
    {á}{{\'a}}1 {é}{{\'e}}1 {í}{{\'i}}1 {ó}{{\'o}}1 {ú}{{\'u}}1
    {Á}{{\'A}}1 {É}{{\'E}}1 {Í}{{\'I}}1 {Ó}{{\'O}}1 {Ú}{{\'U}}1
    {ñ}{{\~n}}1 {Ñ}{{\~N}}1 {ü}{{\"u}}1 {Ü}{{\"U}}1
    {¿}{{?`}}1 {¡}{{!`}}1 {°}{{\textdegree}}1 {º}{{\textordmasculine}}1 {ª}{{\textordfeminine}}1
    {²}{{\textsuperscript{2}}}1 {³}{{\textsuperscript{3}}}1
    {≥}{{$\geq$}}1 {≤}{{$\leq$}}1 {×}{{$\times$}}1 {→}{{$\rightarrow$}}1 {−}{{-}}1
    {–}{{--}}1 {—}{{---}}1 {·}{{$\cdot$}}1 {…}{{...}}1 {“}{{``}}1 {”}{{''}}1
    {‘}{{`}}1 {’}{{'}}1 {≈}{{$\approx$}}1 {±}{{$\pm$}}1 {α}{{$\alpha$}}1 {χ}{{$\chi$}}1
    {η}{{$\eta$}}1 {λ}{{$\lambda$}}1 {∑}{{$\sum$}}1 {€}{{EUR}}1
}

\usepackage[hidelinks]{hyperref}

% --- Datos de la portada: COMPLETAR ---
\newcommand{\titulo}{Patrones multivariados del tráfico de red y detección de ataques en UNSW-NB15}
\newcommand{\autores}{[Integrante 1]\par [Integrante 2]\par [Integrante 3]\par [Integrante 4]}
\newcommand{\afiliacion}{Universidad Católica de Córdoba\par Facultad de [completar]\par Carrera de [completar]\par Sede Río Cuarto}
\newcommand{\materia}{Estadística Multivariada}
\newcommand{\docente}{Prof. Silvana Malpassi}
\newcommand{\fechaentrega}{[día] de [mes] de 2026}

\begin{document}

% ===================== PORTADA =====================
% (no se usa titlepage porque reinicia el contador de páginas)
\begingroup
  \thispagestyle{plain}
  \centering
  \vspace*{2.5cm}
  {\Large\bfseries \titulo \par}
  \vspace{1.2cm}
  {\large Proyecto integrador \par}
  \vspace{2cm}
  {\large \autores \par}
  \vspace{1.5cm}
  {\afiliacion \par}
  \vspace{1.5cm}
  {\textbf{Materia:} \materia \par}
  {\textbf{Docente a cargo:} \docente \par}
  \vspace{1.5cm}
  {\textbf{Fecha de entrega:} \fechaentrega \par}
  \vfill
\endgroup
\clearpage

% ===================== ÍNDICES =====================
\begin{singlespace}
\tableofcontents
\clearpage
\listoffigures
\vspace{1cm}
\listoftables
\end{singlespace}
\clearpage

% ===================== CUERPO =====================
\label{inicio-cuerpo}
\section{Introducción}
\label{sec:introduccion}

[Pendiente]


\section{Datos}
\label{sec:datos}

\textbf{Origen y unidad de análisis.} Trabajamos con UNSW-NB15 \citep{moustafa2015,moustafa2016}, publicada por UNSW Canberra para uso académico con obligación de citar a sus autores. El tráfico normal es realista, pero los ataques se generaron de forma sintética en un laboratorio. La unidad de análisis es el flujo; el prefijo \texttt{s} indica el sentido origen$\to$destino y \texttt{d} el inverso. Partimos del archivo de entrenamiento (175.341 flujos y 45 columnas) y construimos una muestra de 5.000 flujos únicos, estratificada de forma proporcional por \texttt{attack\_cat}: 48,2~\% normales y 51,8~\% ataques (Exploits 18,4~\%, Fuzzers 15,0~\%, Reconnaissance 7,0~\% y el resto, 11,5~\%, agrupado como ``otros ataques''). Al eliminar los duplicados, que en el archivo original eran el 38,6~\% y en su mayoría ataques automáticos, analizamos perfiles distintos de flujo: las proporciones de la muestra no son prevalencias reales de la red.

\textbf{Transformación.} Las variables son positivas, con colas muy largas y rangos de varios órdenes de magnitud. Aplicamos $\log(x+1)$, que conserva los ceros y reduce la asimetría de \texttt{sbytes} de 33,3 a 1,0, la de \texttt{spkts} de 30,9 a 0,8 y la de \texttt{dbytes} de 36,2 a $-0,4$. Persisten colas en \texttt{log\_dur} (2,8), \texttt{log\_ct\_src\_dport\_ltm} (3,1) y \texttt{log\_ct\_dst\_sport\_ltm} (4,7). Los gráficos antes y después están en el Anexo A (Figuras~\ref{fig:a-hist-crudos} y~\ref{fig:a-hist-log}).

\textbf{Variables.} El conjunto X\_num queda formado por las 12 variables de la Tabla~\ref{tab:xnum}, todas en escala logarítmica, agrupadas en los tres bloques esperados. El conjunto X\_cat tiene siete variables: \texttt{tipo\_trafico} (Normal, Exploits, Fuzzers, Reconnaissance y Otros ataques), \texttt{proto\_g} (tcp 70,8~\%, udp 21,9~\%, otros 7,3~\%), \texttt{service\_g} (ninguno 60,0~\%, http 17,4~\%, otros 12,3~\%, dns 10,3~\%), \texttt{state\_g} (FIN 69,3~\%, INT 18,3~\%, CON\_otros 12,4~\%) y tres numéricas discretizadas, \texttt{respuesta\_destino}, \texttt{duracion\_banda} y \texttt{bytes\_origen\_banda}, cuyos cortes se justifican en la sección~\ref{sec:acm}. Ninguna categoría baja del 5~\% (la menor es Reconnaissance, con 7,0~\%). Las variables \texttt{tipo\_trafico}, \texttt{attack\_cat} y \texttt{label} no intervienen en la construcción de componentes, factores ni clusters: solo colorean o se cruzan después.

\begin{table}[H]
  \centering
  \caption{Variables de X\_num (todas en escala $\log(x+1)$; la unidad es la de la variable original).}\label{tab:xnum}
  \begin{tablacompacta}
  \begin{tabular}{@{}llp{6.6cm}l@{}}
  \toprule
  Variable & Bloque esperado & Significado & Unidad \\
  \midrule
  \texttt{spkts} & Volumen & Paquetes enviados por el origen & paquetes \\
  \texttt{sbytes} & Volumen & Bytes enviados por el origen & bytes \\
  \texttt{dbytes} & Volumen & Bytes enviados por el destino & bytes \\
  \texttt{dur} & Temporización & Duración total del flujo & segundos \\
  \texttt{sinpkt} & Temporización & Tiempo entre paquetes del origen & ms \\
  \texttt{sjit} & Temporización & Jitter (variabilidad de esos tiempos) del origen & ms \\
  \texttt{synack} & Temporización & Tiempo entre SYN y SYN-ACK (0 si no es TCP) & segundos \\
  \texttt{ct\_srv\_src} & Repetición & Últimas 100 conexiones con igual servicio y origen & conexiones \\
  \texttt{ct\_srv\_dst} & Repetición & Ídem con igual servicio y destino & conexiones \\
  \texttt{ct\_src\_ltm} & Repetición & Ídem con igual origen & conexiones \\
  \texttt{ct\_src\_dport\_ltm} & Repetición & Ídem con igual origen y puerto destino & conexiones \\
  \texttt{ct\_dst\_sport\_ltm} & Repetición & Ídem con igual destino y puerto origen & conexiones \\
  \bottomrule
  \end{tabular}
  \end{tablacompacta}
\end{table}

\textbf{Calidad.} No hay valores faltantes ni textos vacíos, filas duplicadas ni identificadores repetidos, y tampoco valores negativos, TTL fuera de 0--255, conteos no enteros ni flujos con menos bytes que paquetes (Anexo A, Tabla~\ref{tab:a-calidad}). La única incidencia es un flujo con \texttt{is\_ftp\_login} = 4, un código no binario; lo conservamos porque esa variable no entra en los análisis. El servicio ``\texttt{-}'' se recodificó como ``ninguno'' (sin servicio identificado, no faltante). Hay ceros estructurales que no imputamos: \texttt{synack} en 1.482 flujos (29,6~\%, sin handshake TCP), \texttt{sjit} en 1.294 (25,9~\%) y \texttt{dbytes} en 945 (18,9~\%, el destino no respondió). Para los atípicos usamos el criterio de Tukey, fuera de $[Q_1-1{,}5\,\mathrm{RIC};\,Q_3+1{,}5\,\mathrm{RIC}]$ sobre la escala logarítmica. Resultan 1.714 flujos (34,3~\%) con al menos una de las 12 variables fuera de los límites; la mayor cifra es la de \texttt{log\_dbytes}, 1.056 casos (21,1~\%), de los cuales 945 son los ceros de flujos sin respuesta. En \texttt{log\_ct\_dst\_sport\_ltm} el criterio es poco informativo, porque el 90,7~\% de los flujos vale 1 y el rango intercuartílico es nulo. Decidimos conservar todos los flujos: en ciberseguridad lo extremo puede ser el fenómeno de interés y un valor extremo no prueba un error (boxplots en Anexo A, Figura~\ref{fig:a-boxplots}).

\textbf{Descriptiva y supuestos.} La Figura~\ref{fig:correlaciones} muestra tres zonas. Las variables de volumen se correlacionan entre sí (\texttt{spkts}--\texttt{sbytes} 0,86; \texttt{spkts}--\texttt{dbytes} 0,80), las de temporización también (\texttt{sinpkt}--\texttt{sjit} 0,85; \texttt{dur}--\texttt{sinpkt} 0,68) y las de repetición forman un tercer grupo (\texttt{ct\_srv\_src}--\texttt{ct\_srv\_dst} 0,85; \texttt{ct\_src\_dport\_ltm}--\texttt{ct\_dst\_sport\_ltm} 0,74). Las repeticiones se correlacionan de forma negativa con el resto (entre $-0,13$ y $-0,43$), y \texttt{sjit} tiene correlaciones de 0,46 a 0,55 con las de volumen, de modo que los bloques no están completamente aislados.

\begin{figure}[H]
  \centering
  \grafico[0.56\linewidth]{04_correlaciones.png}
  \caption{Correlaciones de Pearson entre las 12 variables de X\_num. Se distinguen los bloques de volumen, temporización y repetición de conexiones; el último se asocia de forma negativa con los otros dos. Fuente: elaboración propia con UNSW-NB15.}
  \label{fig:correlaciones}
\end{figure}

Los supuestos se cumplen: el determinante de la matriz de correlaciones es $6{,}3\times10^{-5}$ (mayor que $10^{-5}$), el KMO global es 0,748, la menor MSA es 0,621 (\texttt{log\_spkts}; ninguna variable queda bajo 0,50, Anexo A, Figura~\ref{fig:a-msa}) y la prueba de Bartlett rechaza que la matriz sea identidad ($\chi^2$ = 48.303,5; 66 gl; p < 0,001), aunque con n = 5.000 esto ocurre ante casi cualquier asociación y pesa más el KMO. Con las 14 variables candidatas el determinante era $1{,}0\times10^{-7}$, por la redundancia de \texttt{dpkts} (r = 0,95 con \texttt{dbytes}) y de \texttt{rate} (r = $-0,91$ con \texttt{sinpkt}, y calculada a partir de los paquetes y la duración). Quitar solo \texttt{dpkts} dejaba el determinante en $5{,}2\times10^{-6}$, aún bajo el umbral; por eso excluimos las dos. Estas columnas siguen en la base: \texttt{dpkts} define \texttt{respuesta\_destino}.

\textbf{Vínculo con H1.} Las correlaciones son compatibles con tres bloques, aunque la separación entre volumen y temporización es parcial; las secciones siguientes ponen a prueba si ACP, AF y el cluster de variables los recuperan.

\section{Análisis de componentes principales}
\label{sec:acp}

\textbf{Estandarización.} Aplicamos \texttt{prcomp(X\_num, scale. = TRUE)}. Las varianzas de las variables en escala logarítmica son muy desiguales: van de 0,002 (\texttt{log\_synack}) a 13,61 (\texttt{log\_sjit}), más de 7.500 veces, y tres de ellas (\texttt{log\_sjit}, \texttt{log\_dbytes} y \texttt{log\_sinpkt}) reúnen el 83,2~\% de la varianza total. Sin estandarizar, la primera componente captaría el 66,4~\% de la variabilidad, frente al 43,3~\% con estandarización: ese resultado reflejaría las unidades y no las relaciones entre variables. Por eso trabajamos con la matriz de correlaciones \citep{johnson2007}.

\textbf{Cantidad de componentes.} Los autovalores son 5,195, 2,026 y 1,721 para las tres primeras componentes, y el cuarto cae a 0,759; solo tres superan 1 (criterio de Kaiser). El gráfico del codo (Figura~\ref{fig:acp-biplot}, izquierda) muestra una caída marcada hasta la tercera y un tramo casi plano después. Las tres componentes explican el 43,3~\%, el 16,9~\% y el 14,3~\% de la varianza, 74,5~\% acumulado, y cada una de las restantes aporta menos del 6,4~\%. Retenemos tres, que son interpretables y coinciden en número con los bloques de H1.

\textbf{Interpretación.} La Tabla~\ref{tab:acp-cargas} muestra las correlaciones variable--componente. En CP1, las siete variables de volumen y temporización tienen signo positivo (0,54 a 0,82) y las cinco de repetición, negativo ($-0,56$ a $-0,70$); las que más contribuyen son \texttt{log\_sjit} (13,1~\%), \texttt{log\_sinpkt} (9,9~\%) y \texttt{log\_ct\_srv\_dst} (9,4~\%). CP2 se construye sobre todo con \texttt{log\_spkts} (20,1~\%), \texttt{log\_sbytes} (15,7~\%) y las repeticiones \texttt{log\_ct\_srv\_src} (12,7~\%) y \texttt{log\_ct\_srv\_dst} (11,8~\%), todas con signo positivo. CP3 se define por \texttt{log\_sinpkt} (20,1~\%) y \texttt{log\_synack} (19,6~\%), con signo contrario al de los bytes y paquetes. La variable peor representada es \texttt{log\_ct\_src\_ltm} (cos2 total 0,52); la mejor, \texttt{log\_spkts} (0,94).

\begin{table}[H]
  \centering
  \caption{Correlaciones entre cada variable y las tres componentes retenidas, en negrita las de $|r| \geq 0{,}5$, y calidad de representación acumulada (suma de cos2 en CP1--CP3).}\label{tab:acp-cargas}
  \begin{tablacompacta}
  \begin{tabular}{@{}llrrrr@{}}
  \toprule
  Variable & Bloque esperado & CP1 & CP2 & CP3 & cos2 total \\
  \midrule
  \texttt{log\_spkts} & Volumen & \textbf{0,64} & \textbf{0,64} & $-$0,35 & 0,94 \\
  \texttt{log\_sbytes} & Volumen & \textbf{0,60} & \textbf{0,56} & $-$0,32 & 0,78 \\
  \texttt{log\_dbytes} & Volumen & \textbf{0,64} & 0,45 & $-$0,39 & 0,76 \\
  \texttt{log\_dur} & Temporización & \textbf{0,64} & 0,27 & 0,40 & 0,64 \\
  \texttt{log\_sinpkt} & Temporización & \textbf{0,72} & 0,03 & \textbf{0,59} & 0,86 \\
  \texttt{log\_sjit} & Temporización & \textbf{0,82} & 0,23 & 0,36 & 0,87 \\
  \texttt{log\_synack} & Temporización & \textbf{0,54} & 0,00 & \textbf{0,58} & 0,63 \\
  \texttt{log\_ct\_srv\_src} & Repetición & \textbf{$-$0,68} & \textbf{0,51} & 0,08 & 0,73 \\
  \texttt{log\_ct\_srv\_dst} & Repetición & \textbf{$-$0,70} & 0,49 & 0,04 & 0,73 \\
  \texttt{log\_ct\_src\_ltm} & Repetición & \textbf{$-$0,56} & 0,45 & 0,00 & 0,52 \\
  \texttt{log\_ct\_src\_dport\_ltm} & Repetición & \textbf{$-$0,69} & 0,43 & 0,41 & 0,82 \\
  \texttt{log\_ct\_dst\_sport\_ltm} & Repetición & \textbf{$-$0,61} & 0,29 & 0,44 & 0,65 \\
  \bottomrule
  \end{tabular}
  \end{tablacompacta}
\end{table}

\textbf{Nombre de las componentes.} CP1 opone el tamaño y el ritmo del flujo (más paquetes y bytes, más duración, mayor espaciado y jitter) a la repetición de conexiones: la llamamos \emph{magnitud del flujo frente a repetición}. Un puntaje alto describe flujos largos o voluminosos, y uno bajo, ráfagas breves entre los mismos equipos. CP2 es ortogonal a la anterior y reúne volumen (paquetes y bytes) con repetición, ambos en positivo; la llamamos \emph{volumen de tráfico}, aunque no es un eje puro, porque las repeticiones también cargan sobre ella (0,29 a 0,51). CP3 queda definida por el espaciado entre paquetes y el tiempo del handshake, opuestos a los bytes: \emph{ritmo lento}.

\begin{figure}[H]
  \centering
  \grafico[0.9\linewidth]{08_acp_codo_biplot.png}
  \caption{Gráfico del codo (izquierda) y biplot de CP1 y CP2 (derecha), con los flujos coloreados por \texttt{tipo\_trafico}. Los tipos de tráfico se solapan en gran medida; los tres flujos extremos están rotulados con su identificador. Fuente: elaboración propia con UNSW-NB15.}
  \label{fig:acp-biplot}
\end{figure}

\textbf{Observaciones extremas.} Medimos la distancia de cada flujo al centro en el espacio de las tres componentes, ponderada por la varianza de cada una. Los tres más alejados (id 127446, 125215 y 136451, con distancias de 43,6, 38,0 y 32,6) tienen el mismo perfil: tráfico UDP al servicio DNS en estado INT, de 2 paquetes y 114 bytes, sin respuesta y con duración de 6 a 9 microsegundos. Son del tipo ``otros ataques'' (categoría Generic). Lo que los vuelve extremos es la repetición: sus conteos \texttt{ct\_srv\_src} van de 45 a 50 (la mediana es 3) y \texttt{ct\_src\_ltm} de 33 a 50 (mediana 2), lo que explica el puntaje muy negativo en CP1 ($-8,7$ a $-9,6$). No son errores: 152 flujos de la muestra comparten ese patrón de tamaño y todos son ataques. En el biplot, el tipo de tráfico no separa grupos nítidos.

\textbf{Vínculo con H1.} El ACP necesita tres componentes, como bloques se esperaban, y aísla el bloque de repetición, que se opone al resto en CP1. En cambio, volumen y temporización no se separan en CP1 (ambos con signo positivo) y \texttt{log\_sjit} carga más en CP1 que en CP3. Hasta aquí H1 se sostiene en parte; la rotación del análisis factorial dirá si los tres bloques emergen de forma más limpia.

\section{Análisis factorial}
\label{sec:af}

[Pendiente]


\section{Análisis de conglomerados}
\label{sec:cluster}

[Pendiente]


\section{Análisis de correspondencias simple y múltiples}
\label{sec:acm}

\subsection{Construcción de X\_cat}
X\_cat tiene siete variables (Tabla~\ref{tab:xcat}). \texttt{tipo\_trafico} entra como variable activa, a diferencia del ACP, el AF y los conglomerados, porque se piden asociaciones entre categóricas, incluida la de contexto; son descriptivas y en parte esperables. Todas las categorías superan el 5\,\% (la menor, 6,98\,\%), por lo que no fusionamos ninguna. Para \texttt{dur} y \texttt{sbytes}, muy asimétricas (\texttt{sbytes}: mediana 918 y media 16.918 bytes), usamos terciles: cortes de igual amplitud dejarían casi todos los flujos en una categoría, y los terciles dan 33\,\% por categoría sin usar la etiqueta. \texttt{respuesta\_destino} separa los flujos con \texttt{dpkts} = 0 (``sin respuesta''), un corte natural.

\begin{table}[htbp]
  \centering
  \caption[Variables y categorías finales de X\_cat]{X\_cat: variables, origen y categorías finales (\% de los 5.000 flujos). Fuente: elaboración propia con UNSW-NB15.}\label{tab:xcat}
  \begin{tablacompacta}
  \footnotesize\setlength{\tabcolsep}{3pt}
  \begin{tabularx}{\linewidth}{@{}>{\raggedright\arraybackslash}p{3.3cm}>{\raggedright\arraybackslash}p{3.6cm}X@{}}
    \toprule
    Variable & Origen & Categorías (\%) \\
    \midrule
    \texttt{tipo\_trafico} & \texttt{attack\_cat} agrupada & Normal 48,2 · Exploits 18,4 · Fuzzers 15,0 · Otros ataques 11,5 · Reconnaissance 7,0 \\
    \texttt{proto\_g} & \texttt{proto} agrupada & tcp 70,8 · udp 21,9 · otros 7,3 \\
    \texttt{service\_g} & \texttt{service} agrupada & ninguno 60,0 · http 17,4 · otros 12,3 · dns 10,3 \\
    \texttt{state\_g} & \texttt{state} agrupada & FIN 69,3 · INT 18,3 · CON\_otros 12,4 \\
    \texttt{respuesta\_destino} & \texttt{dpkts} = 0 o > 0 & con\_respuesta 81,1 · sin\_respuesta 18,9 \\
    \texttt{duracion\_banda} & terciles de \texttt{dur} (cortes 0,0185 s y 0,709 s) & corta 33,3 · media 33,3 · larga 33,3 \\
    \texttt{bytes\_origen\_banda} & terciles de \texttt{sbytes} (cortes 564 y 1.580 bytes) & bajos 33,8 · medios 33,6 · altos 32,6 \\
    \bottomrule
  \end{tabularx}
  \end{tablacompacta}
\end{table}

\subsection{Análisis de correspondencias simple: \texttt{tipo\_trafico} $\times$ \texttt{state\_g}}
Elegimos \texttt{state\_g} por encarnar H3 y ser la más asociada al tipo de tráfico (V de Cramér 0,405, frente a 0,320 de \texttt{service\_g} y 0,302 de \texttt{proto\_g}). La asociación es significativa ($\chi^2 = 1.637{,}2$, 8 gl, $p < 0{,}001$) y dos dimensiones recogen el 100\,\% de la inercia (77,9 y 22,1\,\%). La Dim 1 opone INT ($-1{,}009$) a CON\_otros ($+0{,}662$): INT concentra el 62,3\,\% de Otros ataques, el 33,0\,\% de Fuzzers y el 26,1\,\% de Reconnaissance contra el 5,9\,\% de Normal, y CON\_otros al 24,0\,\% de Normal y no más del 3,8\,\% de cada ataque. Exploits no sigue el patrón (90,4\,\% en FIN; Normal, 70,0\,\%). (mapa en \texttt{13\_acs\_mapa.png}).

\subsection{Análisis de correspondencias múltiples}
Con $Q = 7$ y $J = 23$, la inercia total cruda es $J/Q - 1 = 2{,}286$ y el autovalor medio $1/Q = 0{,}143$ (seis lo superan: criterio poco exigente). Los porcentajes crudos son bajos (Dim 1: 27,4\,\%, autovalor 0,627; Dim 2: 14,0\,\%, 0,321; acumulado 41,4\,\%). \texttt{mjca} del paquete \texttt{ca} \citep{nenadic2007} usa por defecto la corrección de Greenacre (\texttt{lambda = \textquotedbl{}adjusted\textquotedbl{}}; \citealp{greenacre2017}), que reportamos: inercia total ajustada 0,414, Dim 1 = 77,0\,\%, Dim 2 = 10,4\,\%, Dim 3 = 1,8\,\% (gráfico de inercia en \texttt{14\_acm\_inercia.png}); el plano suma 87,4\,\% de inercia ajustada, cifra optimista que debe leerse junto con el 41,4\,\% crudo. Retenemos \textbf{dos dimensiones}: el codo de los autovalores se marca tras la Dim 2 y la tercera aporta menos del 2\,\% ajustado. Contribuciones, cos2 y $\eta^2$ salen de \texttt{FactoMineR::MCA} \citep{le2008,husson2017}.

\begin{figure}[htbp]
  \centering
  \grafico[0.68\linewidth]{15_acm_categorias.png}
  \caption[Mapa de categorías del ACM]{Mapa de categorías del ACM (Dim 1 y Dim 2, porcentajes crudos). A la derecha, flujos breves, sin respuesta y no TCP; a la izquierda, TCP, FIN, mayor duración y volumen. Fuente: elaboración propia con UNSW-NB15.}
  \label{fig:acm-mapa}
\end{figure}

\begin{table}[htbp]
  \centering
  \caption[Contribuciones y cos2 de las categorías del ACM]{Categorías que más contribuyen (Ctr, \%) a las dimensiones 1 y 2, con su cos2 en cada una. Fuente: elaboración propia con UNSW-NB15.}\label{tab:acm-contrib}
  \begin{tablacompacta}
  \footnotesize\setlength{\tabcolsep}{4pt}
  \begin{tabular}{@{}lrrrr@{\hspace{0.8cm}}lrrrr@{}}
    \toprule
    Categoría & Ctr 1 & cos2 1 & Ctr 2 & cos2 2 & Categoría & Ctr 1 & cos2 1 & Ctr 2 & cos2 2 \\
    \midrule
    \texttt{state\_g}: INT & 12,1 & 0,651 & 9,2 & 0,252 & \texttt{state\_g}: FIN & 6,1 & 0,866 & 0,2 & 0,016 \\
    \texttt{respuesta}: sin & 12,1 & 0,656 & 8,7 & 0,242 & \texttt{proto\_g}: tcp & 5,9 & 0,886 & 0,1 & 0,010 \\
    \texttt{duracion}: corta & 11,0 & 0,726 & 0,5 & 0,018 & \texttt{service\_g}: dns & 5,7 & 0,277 & 15,2 & 0,381 \\
    \texttt{proto\_g}: udp & 9,7 & 0,545 & 6,9 & 0,199 & \texttt{state\_g}: CON\_otros & 2,5 & 0,127 & 22,8 & 0,585 \\
    \texttt{bytes\_origen}: bajos & 7,8 & 0,520 & 0,4 & 0,015 & \texttt{proto\_g}: otros & 4,7 & 0,223 & 12,0 & 0,291 \\
    \bottomrule
  \end{tabular}
  \end{tablacompacta}
\end{table}

\textbf{Dim 1: ``conexión efímera sin respuesta'' frente a ``sesión TCP completa''.} La construyen INT, sin\_respuesta, duración corta, UDP y bytes bajos contra FIN y TCP (Tabla~\ref{tab:acm-contrib}); discriminan sobre todo \texttt{state\_g} ($\eta^2 = 0{,}909$), \texttt{proto\_g} (0,891) y \texttt{duracion\_banda} (0,726), y \texttt{tipo\_trafico} apenas (0,242; Figura~\ref{fig:acm-clusters}, derecha). \textbf{Dim 2: ``DNS/UDP con estado CON'' frente a ``protocolo \emph{otros} no establecido''.} La construyen CON\_otros (22,8~\%), dns (15,2~\%) y proto \emph{otros} (12,0~\%), y Normal queda del lado DNS/CON (cos2 0,384); los demás tipos tienen cos2 de 0,03 a 0,08 (posición solo indicativa). \textbf{Redundancia:} \texttt{proto\_g} = \emph{otros} implica \texttt{service\_g} = ninguno (363 de 363) y casi siempre INT (344 de 363); FIN solo ocurre en TCP; todo flujo INT es sin\_respuesta (913 de 913) y casi todo flujo sin respuesta es INT (913 de 945). Por eso se agrupan en el mapa.

\begin{figure}[htbp]
  \centering
  \grafico[0.95\linewidth]{16_acm_clusters_eta.png}
  \caption[Clusters suplementarios y eta$^2$ del ACM]{Izquierda: clusters de la Etapa 3 proyectados como suplementarios (rombos) sobre las observaciones; los clusters 1, 2 y 4 se superponen (Dim 1 entre $-0{,}62$ y $-0{,}55$). Derecha: $\eta^2$ por variable en Dim 1 y Dim 2; \texttt{state\_g} y \texttt{proto\_g} dominan la Dim 1. Fuente: elaboración propia con UNSW-NB15.}
  \label{fig:acm-clusters}
\end{figure}

\textbf{Clusters suplementarios} (Figura~\ref{fig:acm-clusters}). No ocupan zonas totalmente distintas: los clusters 1, 2 y 4 se superponen en la zona TCP/FIN y se distinguen por duración y bytes. El cluster 3 (Dim 1 = 2,16) queda aislado a la derecha: 100\,\% INT, sin respuesta, duración corta y bytes bajos, y 87,8\,\% UDP y dns. El 5 (Dim 1 = 1,37) es breve (96,1\,\% corta), UDP (73,9\,\%) y con 45,6\,\% CON\_otros; el 6 es intermedio. Los cos2 en el plano son bajos salvo el 5 (0,59): lectura parcial.

\textbf{Tres asociaciones.} (i) \emph{Otros ataques con INT, sin respuesta y flujos pequeños}: 62,3\,\% INT y 63,2\,\% sin respuesta (Normal: 5,9 y 6,9\,\%), 62,8\,\% con bytes bajos. (ii) \emph{Normal con CON\_otros y con respuesta} (24,0\,\% y 93,1\,\%); el DNS no es exclusivo de lo normal (14,4\,\% frente a 27,6\,\% en Otros ataques). (iii) \emph{Exploits con TCP/FIN, duración larga y http}: 54,7\,\% larga (Normal: 23,7\,\%) y 44,2\,\% http (Normal: 9,9\,\%). Son frecuencias de laboratorio; el ACM es descriptivo.

\section{Regresión logística}
\label{sec:logistica}

\textbf{Diseño.} Para contrastar H4 modelamos la probabilidad de que un flujo sea ataque (\texttt{label} = 1) con regresión logística \citep{hosmer2013}. Dividimos la muestra en 3.500 flujos de entrenamiento (51,8\,\% ataques) y 1.500 de prueba (51,9\,\%), estratificando por \texttt{label} (semilla 42) (no usamos el testing set original: no está en el repositorio y habría que prepararlo igual). Para evitar fuga reajustamos el ACP solo con el entrenamiento y proyectamos la prueba; las cargas son casi idénticas a las del capítulo~\ref{sec:acp} (correlación $\ge 0{,}999$). Los predictores son PC1 a PC3 y \texttt{proto\_g} (referencia tcp), \texttt{service\_g} (ninguno) y \texttt{state\_g} (FIN), con la categoría más frecuente como referencia. Excluimos TTL (separación perfecta artefactual) y \texttt{tipo\_trafico} y \texttt{attack\_cat} (fuga de la etiqueta).

\textbf{Elección del modelo.} M1 usa solo los puntajes; M2 agrega protocolo, servicio y estado. La razón de verosimilitud favorece a M2 ($\chi^2 = 862{,}2$; 7 gl; $p < 0{,}001$), con unos 850 puntos menos de AIC y 0,08 más de AUC (Tabla~\ref{tab:log-modelos}). Elegimos M2.

\begin{table}[H]
  \centering
  \caption{Comparación de modelos; el AUC es sobre la muestra de prueba.}\label{tab:log-modelos}
  \begin{tablacompacta}
  \begin{tabular}{lrrr} \toprule
  Modelo & AIC & $R^2$ McFadden & AUC prueba \\ \midrule
  M1: 3 puntajes del ACP & 3.744,1 & 0,229 & 0,794 \\
  M2: M1 + protocolo, servicio, estado & 2.895,9 & 0,407 & 0,878 \\ \bottomrule
  \end{tabular}
  \end{tablacompacta}
\end{table}

\textbf{Odds ratios.} Con lo demás fijo, un desvío más de PC1 (más volumen y temporización, menos repetición) multiplica las chances de ataque por 4,7, y uno más de PC2 (más volumen y repetición a la vez) las reduce un 39\,\% (Tabla~\ref{tab:log-or}). El estado CON\_otros tiene chances mucho menores que FIN (OR 0,05). Los OR de las categóricas son condicionales y poco estables por colinealidad (VIF $\approx 16$ en \texttt{udp} e INT; $\le 3{,}7$ en las componentes), así que no deben leerse de a uno: dns tiene 31\,\% de ataques en el entrenamiento (49\,\% en servicio ninguno) pero OR ajustado 18,7: el signo cambia al controlar por estado y protocolo. El IC de INT es muy ancho: no debe citarse ``93 veces''.

\begin{table}[H]
  \centering
  \caption{Odds ratios de M2 con IC 95\,\% por verosimilitud perfilada (componentes: por desvío del puntaje; categóricas: frente a la referencia). Los términos de protocolo (udp, otros) y servicio otros tienen IC que incluyen 1.}\label{tab:log-or}
  \begin{tablacompacta}
  \begin{tabular}{lrr} \toprule
  Término & OR & IC 95\,\% \\ \midrule
  PC1 & 4,68 & 3,82 -- 5,77 \\
  PC2 & 0,61 & 0,54 -- 0,69 \\
  PC3 & 2,55 & 2,27 -- 2,89 \\
  Estado INT vs FIN & 92,7 & 22,6 -- 361,8 \\
  Estado CON\_otros vs FIN & 0,05 & 0,02 -- 0,13 \\
  Servicio dns vs ninguno & 18,7 & 8,5 -- 43,0 \\
  Servicio http vs ninguno & 1,88 & 1,47 -- 2,42 \\ \bottomrule
  \end{tabular}
  \end{tablacompacta}
\end{table}

\textbf{Desempeño en prueba.} Con umbral 0,5, M2 logra AUC = 0,878 (IC 95\,\% de DeLong con \texttt{pROC} \citep{robin2011}: 0,860 -- 0,896), exactitud 0,817, sensibilidad 0,910, especificidad 0,717 y precisión 0,776. Domina el falso positivo: 204 normales marcados como ataque (28,3\,\% de los normales) frente a 70 ataques no detectados (9,0\,\%; Figura~\ref{fig:log-roc}). En detección de intrusiones el modelo prioriza no dejar pasar ataques a costa de alertas falsas.

\begin{figure}[H]
  \centering
  \grafico[0.6\linewidth]{17_logistica_roc_confusion.png}
  \caption{Curva ROC y matriz de confusión de M2 en la prueba (umbral 0,5). La curva se aleja de la diagonal del azar; el error principal son los normales clasificados como ataque. Fuente: elaboración propia con UNSW-NB15.}
  \label{fig:log-roc}
\end{figure}

\textbf{Diagnóstico.} No hay separación: convergió en 6 iteraciones sin avisos, el coeficiente máximo es 4,5 y ninguna probabilidad ajustada está a menos de $10^{-8}$ de 0 o 1. En influencia, 151 flujos (4,3\,\%) superan $4/n$ en la distancia de Cook, pero el máximo es 0,030, muy lejos de 1. La calibración es razonable, salvo que el decil 10 sobreestima (predicho 0,94, observado 0,86; Anexo A, Figura~\ref{fig:a-log-or-cal}).

\textbf{Cierre.} H4 se confirma: el AUC en prueba es 0,878 y el límite inferior de su IC (0,860) también supera 0,80; con las componentes solas (M1, 0,794) no se cumpliría. La probabilidad estimada es relativa a esta base (52\,\% de ataques, sin duplicados, ataques sintéticos): interpretamos odds ratios, no prevalencias, y no extrapolamos a una red real.

\section{Integración y conclusiones}
\label{sec:integracion}

\begin{table}[H]
  \centering
  \caption{Integración de las técnicas: estructura hallada, protagonistas y relación con las demás. Fuente: elaboración propia con UNSW-NB15.}\label{tab:integracion}
  \begin{tablacompacta}
  \footnotesize
  \setlength{\tabcolsep}{3pt}
  \begin{tabularx}{\linewidth}{@{}>{\raggedright\arraybackslash}p{1.45cm}>{\raggedright\arraybackslash}X>{\raggedright\arraybackslash}X>{\raggedright\arraybackslash}X@{}}
  \toprule
  Técnica & Estructura & Protagonistas & Coincidencias y discrepancias \\
  \midrule
  ACP (§\ref{sec:acp}) & 3 componentes, 74,5~\% de la varianza: CP1 magnitud del flujo frente a repetición (43,3~\%), CP2 volumen de tráfico, CP3 ritmo lento & CP1: \texttt{sjit}, \texttt{sinpkt} frente a las cinco \texttt{ct\_*}; CP2: \texttt{spkts}, \texttt{sbytes} & Aísla la repetición, como el AF y el cluster de variables; pero sus ejes son mixtos y no separan volumen de temporización \\
  AF (§\ref{sec:af}) & 3 factores (GLS, varimax), 68,6~\%: F1 Repetición, F2 Temporización, F3 Volumen & Cargas de 0,54 a 0,91, sin cargas cruzadas ambiguas & Coincide exactamente con el cluster de variables (ARI = 1,000) y con H1; ajuste global flojo (TLI 0,609, RMSEA 0,239) \\
  Cluster de variables (§\ref{sec:cluster}) & 3 bloques (Ward): temporización, volumen, repetición & Las cinco \texttt{ct\_*} juntas; \texttt{sinpkt}--\texttt{sjit} (0,85) & Igual al AF; el ACP mezcla los bloques en cada eje \\
  Cluster de observaciones (§\ref{sec:cluster}) & 6 grupos (k-means), silhouette 0,333, estructura débil & C1 transferencias voluminosas (1,4~\% ataque); C3 ráfagas repetidas sin respuesta (100~\%); C2 y C6 mal separados & Los grupos se ordenan en CP1 y F1 (C3: $-7{,}00$ y $+3{,}78$); en el ACM, C3 queda aislado en Dim 1. Con puntajes del ACP, ARI = 0,708 \\
  ACS y ACM (§\ref{sec:acm}) & ACS: tipo $\times$ estado (V = 0,405, 100~\% de inercia en 2 dim.). ACM: 2 dim., 41,4~\% crudo (87,4~\% ajustado) & Dim 1: INT, sin respuesta, duración corta y UDP frente a FIN y TCP; Dim 2: DNS/CON frente a protocolo \emph{otros} & Dim 1 reproduce el contraste de CP1 y el cluster 3; \texttt{tipo\_trafico} casi no pesa ($\eta^2$ = 0,242) \\
  Logística (§\ref{sec:logistica}) & M2 con PC1--PC3 y categorías: AUC 0,882 en prueba (IC 95~\% 0,864--0,900) & PC1 (OR 5,01), PC3 (2,54), PC2 (0,59); INT, CON\_otros (0,19), dns & Coincide con el ACM en INT y CON\_otros; discrepa en dns (signo opuesto al bruto) y los OR categóricos son inestables \\
  \bottomrule
  \end{tabularx}
  \end{tablacompacta}
\end{table}

\textbf{Hipótesis.}
\begin{singlespace}
\begin{itemize}\setlength{\itemsep}{0pt}\setlength{\topsep}{1pt}\setlength{\parsep}{0pt}
  \item \textbf{H1, confirmada.} El AF y el cluster de variables dan los mismos tres bloques (ARI = 1,000; cargas $\ge 0{,}54$). El ACP pide también tres componentes, pero no los recupera como ejes puros (secciones~\ref{sec:acp} y~\ref{sec:af}).
  \item \textbf{H2, parcialmente confirmada.} El cluster 1 es 98,6~\% normal y el 3 es 100~\% ataque, pero los otros cuatro tienen entre 34,9 y 72,6~\% de ataques. Además, el cluster 3 sí distingue un tipo (130 de 147 flujos son Generic): falla la parte de ``no distinguen tipos'' (sección~\ref{sec:cluster}).
  \item \textbf{H3, parcialmente confirmada.} INT y sin respuesta se asocian con Otros ataques (62,3~\% y 63,2~\% frente a 5,9~\% y 6,9~\% de Normal), pero Exploits es casi todo FIN con respuesta, y dns y http no son exclusivos de lo normal (sección~\ref{sec:acm}).
  \item \textbf{H4, confirmada.} AUC de prueba 0,882, con límite inferior 0,864 $>$ 0,80; solo con las componentes (M1) sería 0,795 (sección~\ref{sec:logistica}).
\end{itemize}
\end{singlespace}

\textbf{Bloques y clusters.} Los bloques coinciden entre AF y cluster de variables, y el ACP los recupera solo en parte: separa la repetición, pero mezcla volumen y temporización en CP1. Los clusters se ordenan sobre todo por CP1 ($\eta^2 = 0{,}86$; CP2 y CP3, 0,65) y por los factores (F2 0,78, F3 0,73, F1 0,72): por ejemplo, el cluster 3 tiene F1 = $+3{,}78$ (repetición) y el 2, F2 = $+0{,}88$ (sesiones largas). Es en parte circular, porque se agrupó con las mismas variables. Las categorías del ACM lo confirman con variables que no entraron al cluster: el cluster 3 es 100~\% INT y sin respuesta (Dim 1 = 2,16) y el 5 es breve y UDP (Dim 1 = 1,37), mientras que los clusters 1, 2 y 4 se superponen en TCP/FIN y solo se distinguen por duración y bytes.

\textbf{Sorpresas.} (i) La repetición de conexiones, que esperábamos asociada a los ataques, no los distingue por sí sola: los clusters 3 y 6 la tienen alta, pero uno es 100~\% ataque y el otro 34,9~\%. (ii) El cluster 3 aísla ataques Generic. (iii) INT es redundante con la falta de respuesta (913 de 913 flujos), así que el ACM y la logística no pueden separarlos. (iv) El caso Heywood: \texttt{log\_spkts} supera comunalidad 1 con casi cualquier estimador y solo GLS lo evita. (v) Los OR categóricos son inestables entre particiones (INT de 93 a 375, udp de 0,42 a 0,14). (vi) El ACP, con ejes mixtos, resume más varianza (74,5~\%) pero el AF, con bloques puros, interpreta mejor.

\textbf{Limitaciones.} Los ataques son sintéticos, de laboratorio; sin duplicados, las proporciones no son prevalencias y la probabilidad logística es relativa a esta base (52~\% de ataques); excluimos TTL por artefacto; los \texttt{ct\_*} usan ventanas de 100 conexiones, así que los flujos no son independientes; con $n = 5.000$ Bartlett rechaza cualquier asociación y pesa más el KMO; persiste asimetría pese al logaritmo; el ajuste global del AF es flojo; k = 6 tiene una silhouette casi empatada con k = 2 a 5; los terciles del ACM son un criterio estadístico y no sustantivo; hay colinealidad entre categóricas en la logística (VIF $\approx$ 17--18); usamos una sola partición 70/30; y los datos son observacionales, por lo que describimos asociaciones y no causas.

\textbf{Próximos pasos.} Validar con el testing set oficial procesado igual (mismo logaritmo, cortes y ACP), repetir con validación cruzada, y probar modelos con interacciones o regularización que reduzcan la inestabilidad de los OR.

% REVISAR GRUPO: opinión propia
\textbf{Conclusión.} Respondemos a la pregunta guía así: los flujos se organizan en tres patrones, volumen, temporización y repetición, y esos patrones, junto con el estado, el servicio y el protocolo, permiten estimar la probabilidad de ataque con un AUC de 0,882 en prueba, aunque con 28,7~\% de los normales marcados como ataque. A nuestro juicio (valoración del grupo), lo más sólido es la estructura de variables, que tres técnicas recuperan, y lo más frágil es la lectura de grupos y de odds ratios aislados. Creemos que el estado de la conexión aporta más que cualquier patrón numérico, pero no podemos extrapolar a una red real.

\label{fin-cuerpo}

% ===================== REFERENCIAS (página nueva) =====================
\clearpage
\phantomsection
\addcontentsline{toc}{section}{Referencias bibliográficas}
\renewcommand{\refname}{Referencias bibliográficas}
\bibliographystyle{apalike}
\bibliography{referencias}

% ===================== ANEXOS =====================
\clearpage
\appendix
\renewcommand{\thesection}{\Alph{section}}
% En los anexos, figuras y tablas se numeran A.1, A.2, B.1...
\counterwithin{figure}{section}
\counterwithin{table}{section}
\section{Figuras y tablas complementarias}
\label{anexo:complementos}

[Pendiente]


% Anexo: código fuente. Se lee directamente de ../R y ../python, así que
% el anexo siempre muestra la versión actual de cada script.
\section{Anexo: código fuente}
\label{anexo:scripts}
Los scripts se ejecutan desde la carpeta principal del repositorio. Primero se corre \texttt{R/ejecutar\_todo.R}, que genera las tablas en \texttt{resultados/tablas}; luego los scripts de Python dibujan las figuras a partir de esas tablas. Todos los procesos aleatorios usan la semilla 42.
\subsection*{Scripts de R}
\IfFileExists{../R/00_datos.R}{%
\subsubsection*{\texttt{R/00\_datos.R} · Etapa 0: importación, calidad, descriptiva y supuestos}
\lstinputlisting[language=R]{../R/00_datos.R}
}{}
\IfFileExists{../R/01_acp.R}{%
\subsubsection*{\texttt{R/01\_acp.R} · Etapa 1: análisis de componentes principales}
\lstinputlisting[language=R]{../R/01_acp.R}
}{}
\IfFileExists{../R/02_af.R}{%
\subsubsection*{\texttt{R/02\_af.R} · Etapa 2: análisis factorial}
\lstinputlisting[language=R]{../R/02_af.R}
}{}
\IfFileExists{../R/aux_diagnostico_af.R}{%
\subsubsection*{\texttt{R/aux\_diagnostico\_af.R} · Etapa 2 (auxiliar): diagnóstico del caso Heywood}
\lstinputlisting[language=R]{../R/aux_diagnostico_af.R}
}{}
\IfFileExists{../R/03_cluster.R}{%
\subsubsection*{\texttt{R/03\_cluster.R} · Etapa 3: conglomerados de observaciones y de variables}
\lstinputlisting[language=R]{../R/03_cluster.R}
}{}
\IfFileExists{../R/04_acm.R}{%
\subsubsection*{\texttt{R/04\_acm.R} · Etapa 4: correspondencias simple y múltiples}
\lstinputlisting[language=R]{../R/04_acm.R}
}{}
\IfFileExists{../R/05_integracion.R}{%
\subsubsection*{\texttt{R/05\_integracion.R} · Etapa 6: integración de técnicas}
\lstinputlisting[language=R]{../R/05_integracion.R}
}{}
\IfFileExists{../R/06_logistica.R}{%
\subsubsection*{\texttt{R/06\_logistica.R} · Regresión logística}
\lstinputlisting[language=R]{../R/06_logistica.R}
}{}
\IfFileExists{../R/ejecutar_todo.R}{%
\subsubsection*{\texttt{R/ejecutar\_todo.R} · Ejecución completa del proyecto}
\lstinputlisting[language=R]{../R/ejecutar_todo.R}
}{}
\subsection*{Scripts de Python}
\IfFileExists{../python/graficos_etapa0.py}{%
\subsubsection*{\texttt{python/graficos\_etapa0.py} · Gráficos de la Etapa 0}
\lstinputlisting[language=Python]{../python/graficos_etapa0.py}
}{}
\IfFileExists{../python/graficos_parte2.py}{%
\subsubsection*{\texttt{python/graficos\_parte2.py} · Gráficos de las Etapas 1 a 4}
\lstinputlisting[language=Python]{../python/graficos_parte2.py}
}{}
\IfFileExists{../python/graficos_logistica.py}{%
\subsubsection*{\texttt{python/graficos\_logistica.py} · Gráficos de la regresión logística}
\lstinputlisting[language=Python]{../python/graficos_logistica.py}
}{}
\IfFileExists{../python/verificar_resultados.py}{%
\subsubsection*{\texttt{python/verificar\_resultados.py} · Verificación independiente de resultados}
\lstinputlisting[language=Python]{../python/verificar_resultados.py}
}{}


\end{document}
